% ID: SEL-D0071-OA-E02
% ESTADO: validada. S071.
\textbf{a)} El corte con el eje horizontal se obtiene de $x^3-1=0$ y es $(1,0)$. El corte con el eje vertical es $(0,-1)$. La función es estrictamente creciente en todo $\mathbb R$, pues $x^3$ lo es. Además, $f'(x)=3x^2$ es positiva salvo en $0$. Allí no cambia de signo, por lo que no hay extremos relativos.

\textbf{b)} $f''(x)=6x$. La gráfica es cóncava hacia abajo en $(-\infty,0)$ y hacia arriba en $(0,\infty)$. El cambio de curvatura ocurre en el punto de inflexión $\boxed{(0,-1)}$.

\textbf{c)} $3x^2=3$ da $x=-1$ o $x=1$. Los puntos son $\boxed{(-1,-2)\text{ y }(1,0)}$.
